\section{A current on the complete original physical source}
\label{sec:v72-complete-current}
A representation on the selected disks alone leaves the first-incidence
and outside sources untouched. We instead differentiate their positive
sum on the original collision space. The result has no dependence on
a choice of Morse disks or on an ordering of artificial chart cuts.

Fix $m,R,\lambda=(n,k)$ and $\varepsilon>0$. Let
$\Pi_{n,k,m,R}$ be the original exact-return event. Its occupation
count uses collisions $0,\ldots,m-1$, and its terminal section test
is made at collision $m$. On this event the original roof is
$L_{m,R}=T_{n,R}$. Write
\begin{equation}\label{eq:v72-complete-source-density}
 b=b^\varepsilon_{n,k,m,R}
 =\frac{d\mu^\varepsilon_{n,k,m,R}}{dt},\qquad
 \mu^\varepsilon_{n,k,m,R}=(L_{m,R})_\#
  \bigl((1-H_{m,R}^\varepsilon)\one_{\Pi_{n,k,m,R}}\nu_R^*\bigr).
\end{equation}
Here the first quotient denotes the Radon--Nikodym density of the displayed
measure, not its distributional derivative. The density exists by
the inherited coarea construction. It is nonnegative, integrable
and supported in the finite physical roof range at this count.
The probability $\nu_R^*$ already contains $c_*^{-1}$.

\begin{theorem}[Complete physical current and partition independence]
\label{thm:v72-complete-current}
There is a canonical compactly supported distribution
$\mathfrak D^\varepsilon_{n,k,m,R}=Db$ given by
\begin{equation}\label{eq:v72-physical-current-definition}
 \langle\mathfrak D^\varepsilon_{n,k,m,R},\phi\rangle
 =-\int_{\Pi_{n,k,m,R}}(1-H_{m,R}^\varepsilon(x))
                    \phi'(L_{m,R}(x))\,d\nu_R^*(x).
\end{equation}
It has order at most one. On every common source disintegration,
\begin{equation}\label{eq:v72-current-partition}
 \begin{split}
 \mathfrak D^\varepsilon
 &=D b^{\varepsilon,\mathrm{inc}}
       +D b^{\varepsilon,\mathrm{clr}}\\
 &=D s_{71,1}+D d_{71,1}
       +D b^{\varepsilon,\mathrm{inc}}+D e_{70}.
 \end{split}
\end{equation}
These are distributional identities on the whole roof line. In
particular artificial disk boundaries and first-defect subdivisions
cannot create an additional current in their sum. Physical itinerary
boundaries and boundaries of the exact-label event are not erased.
\end{theorem}
\begin{proof}
The right side of \eqref{eq:v72-physical-current-definition} is
bounded in absolute value by $\|\phi'\|_\infty$ times the mass
of the displayed positive source. Pushforward gives
$-\int b(t)\phi'(t)\,dt$, which is precisely $Db$. This proves
existence, the support assertion, and the order bound. The two source
partitions are \eqref{eq:v51-exact-positive-split} and
\eqref{eq:v71-complete-source}, respectively, with the unchanged
incidence source added. Differentiate their exact equalities. The
result depends only on the original source in
\eqref{eq:v72-complete-source-density}, not on any partition used
to compute it. For bounded measurable insertions the same construction
holds by linearity and domination in variation; no regularity of
such an insertion is inferred.
\end{proof}

We record how the distribution is computed geometrically without
continuing an orbit through a singularity. This also fixes the
orientation of its boundary terms.

\begin{proposition}[Regular coarea flux and exhaustion]
\label{prop:v72-coarea-flux}
In a regular coordinate patch of the original exact-label event,
let $F=L_{m,R}$, let $\Omega$ be a bounded $C^1$ subdomain with
compact closure and $|\nabla F|>0$, and let $w$ be its $C^1$
physical source density. Put $V=\nabla F/|\nabla F|^2$ and let
$b_\Omega$ be the density of $F_\#(w\one_\Omega\,dx)$. Then
\begin{equation}\label{eq:v72-coarea-flux}
 Db_\Omega=F_\#\bigl(\operatorname{div}(wV)\,dx|_\Omega
                   -(wV\cdot n)\,d\mathcal H^1|_{\partial\Omega}\bigr).
\end{equation}
The same formula holds piecewise after common interfaces are combined
with their orientations. Regular compact exhaustions of the original
source converge to \eqref{eq:v72-physical-current-definition} in
distributions. This passage requires only convergence of source mass;
it does not assert convergence in total variation of the currents.
\end{proposition}
\begin{proof}
Since $V\cdot\nabla F=1$, the divergence theorem applied to
$\phi(F)wV$ gives
\[
 \int_\Omega\phi'(F)w\,dx
 =\int_{\partial\Omega}\phi(F)wV\cdot n\,d\mathcal H^1
                  -\int_\Omega\phi(F)\operatorname{div}(wV)\,dx.
\]
Its negative is the action of $Db_\Omega$, proving the formula.
On a common interface carrying the same physical roof and density,
the two outward normals have opposite signs. The interface terms
therefore cancel. If the physical density has a true jump, their sum
is instead the jump flux and must be retained. This assertion never
identifies roofs or sources belonging to different exact labels.

For the exhaustion, first restrict the physical measure to regular
coordinate patches and remove neighborhoods of the critical set.
The singular set and the critical set have collision measure zero
in the inherited finite-word coarea construction. On the remaining
open cells, regular compact subdomains and smooth density cutoffs
approximate the source in $L^1$. Piecewise smooth guards may be
approximated there in the same way. All cutoffs are applied to actual
physical states, not to their nonphysical continuations. If $b_q$
is the resulting roof density, positive domination, or the contraction
of variation under pushforward, gives $\|b_q-b\|_1\to0$.
Consequently
\[
 |\langle Db_q-Db,\phi\rangle|
                  \le\|b_q-b\|_1\|\phi'\|_\infty\longrightarrow0.
\]
This proves the asserted distributional limit and its independence
of the exhaustion. Boundary flux near a critical point is retained
in this limit, even when the corresponding source mass tends to zero.
The entrance atom of a Morse step is the elementary example. Neither
a uniform bound on $\|Db_q\|_{\TV}$ nor a limit exchange with
$m\to\infty$ has been made.
\end{proof}

In particular, for smooth physical factors $a_1,\ldots,a_N$ the
first-defect identity and its derivative are
\[
 \sum_{i=1}^N(1-a_i)\prod_{r<i}a_r=1-\prod_{i=1}^Na_i,
 \qquad
 D\!\left(\sum_{i=1}^N(1-a_i)\prod_{r<i}a_r\right)
                           =-D\!\left(\prod_{i=1}^Na_i\right).
\]
Thus derivatives of an artificial witness subdivision are combined
before a Jordan decomposition. This does not permit cancellation of
positive incidence and clearance densities: it is cancellation in
the derivative of their unchanged positive sum.

\subsection{A fixed-width primitive of the complete current}
For $\delta>0$ set
\begin{equation}\label{eq:v72-average-kernel}
 \rho_\delta(t)=\frac{\one_{(-\delta,\delta)}(t)}{2\delta},
 \qquad A_\delta b=\rho_\delta*b,
 \qquad
 \kappa_\delta(t)=
 \begin{cases}
 -(\delta+t)/(2\delta),&-\delta<t<0,\\
 (\delta-t)/(2\delta),&0<t<\delta,\\
 0,&|t|\ge\delta.
 \end{cases}
\end{equation}
Its value at zero is irrelevant to the distribution. Notice its
unit upward jump there.

\begin{theorem}[Exact current primitive at the original roof]
\label{thm:v72-current-primitive}
Convolution in distributions gives the $L^1$ function
\begin{equation}\label{eq:v72-current-primitive}
 \kappa_\delta*\mathfrak D^\varepsilon=b-A_\delta b.
\end{equation}
Thus the equality holds almost everywhere on the original roof line.
Its Fourier multiplier is
\begin{equation}\label{eq:v72-current-multiplier}
 \widehat\kappa_\delta(\xi)
  =\frac{1-\sin(\delta\xi)/(\delta\xi)}{i\xi}\quad(\xi\ne0),
 \qquad \widehat\kappa_\delta(0)=0.
\end{equation}
In particular
$|\widehat\kappa_\delta(\xi)|
 \le\min\{\delta^2|\xi|/6,2/|\xi|\}$ for $\xi\ne0$.
\end{theorem}
\begin{proof}
Direct distributional differentiation gives
$D\kappa_\delta=\delta_0-\rho_\delta\,dt$. The unit atom is
essential. Both $\kappa_\delta$ and $b$ have compact support at
the fixed count, so differentiation commutes with their convolution:
\[
 \kappa_\delta*Db=(D\kappa_\delta)*b=b-\rho_\delta*b.
\]
The right side is in $L^1$ and fixes its almost-everywhere version.
The Fourier transform of the box average is
$\sin(\delta\xi)/(\delta\xi)$. This proves the multiplier formula
away from zero. Its limit at zero is zero. Finally
$1-\sin x/x=\int_0^1(1-\cos(rx))\,dr$ is nonnegative and at
most $x^2/6$, and its absolute value is at most two.
\end{proof}

This identity is valid even if $Db$ is not a finite measure. It can
also be read by first convolving $Db$ with a smooth approximate
identity at each fixed count, and then taking that approximation
parameter to zero. A positive part in the next section is taken
only after \eqref{eq:v72-current-primitive} has produced an actual
function. The factor $1/|\xi|$ alone supplies no integrable
high-frequency majorant and no transfer-operator resolvent estimate.
